% SOURCE_STATEMENT: "Derive Nekrasov's blow-up identity for the $N_f=4$, $\\operatorname{SU}(2)$ instanton partition function from the affine–Virasoro vertex-operator-algebra identity under $\\frac{\\epsilon_1-\\epsilon_2}{\\epsilon_2}=b^2=-\\frac{k+3}{k+2}$, including the required normalization of conformal blocks."
% Preparation and final self-audit: CONTRIBUTING.md prompts C and D; source/audit artifacts retained in the local physics_statements_2026_10_11 batch.
% !TeX program = lualatex
\documentclass[11pt]{article}
\usepackage[margin=25mm]{geometry}
\usepackage{fontspec}
\setmainfont{Latin Modern Roman}
\usepackage{amsmath,amssymb,mathtools}
\usepackage{unicode-math}
\setmathfont{Latin Modern Math}
\usepackage{xurl,hyperref}
\hypersetup{colorlinks=true,urlcolor=blue}
\setcounter{secnumdepth}{0}
\setlength{\parindent}{0pt}
\setlength{\parskip}{5pt}
\setlength{\emergencystretch}{3em}
\newcommand{\sref}[2]{\hyperlink{source:#1:#2}{[S#2]}}
\newcommand{\cataloguescope}[1]{\paragraph{Recorded target.}#1}
\newcommand{\R}{\mathbb R}
\newcommand{\C}{\mathbb C}
\newcommand{\N}{\mathbb N}
\newcommand{\Z}{\mathbb Z}
\newcommand{\Q}{\mathbb Q}
\newcommand{\F}{\mathbb F}
\newcommand{\eps}{\varepsilon}
\begin{document}
% UnsolvedMath ID 30005075; source catalogue code OWR-9790367-005
\setcounter{section}{30005074}
\section{Affine–Virasoro Derivation of Nekrasov Blow-Up Identities}\label{30005075}
{\small UnsolvedMath ID 30005075 · Mathematical Physics}
\subsection{Definitions and mathematical statement}

\paragraph{Shared notation.}$\mathbb N=\{1,2,\ldots\}$; the standard number fields have their usual meanings. All additional conventions are stated locally.


Write $\Psi(a,m;w,q;\epsilon_1,\epsilon_2)$ for the instanton partition function with surface defect and $Z(a,m;q;\epsilon_1,\epsilon_2)$ for the function without that defect. Here $a$ is the Coulomb parameter, $m$ the four masses, $w$ the defect parameter, $q$ the gauge coupling, and $\epsilon_1,\epsilon_2$ the deformation parameters. The source blow-up relation reads
\[
\Psi(a,m;w,q;\epsilon_1,\epsilon_2)=\sum_{n\in\mathbb Z}
\Psi(a+\epsilon_1n,m;w,q;\epsilon_1,\epsilon_2-\epsilon_1)
Z(a+\epsilon_2n,m;q;\epsilon_1-\epsilon_2,\epsilon_2).
\]
An affine--Virasoro identity decomposes affine Lie-algebra representations into factors involving Virasoro representations.

\paragraph{Catalogue wording retained for traceability.}
Derive Nekrasov's blow-up identity for the $N_f=4$, $\operatorname{SU}(2)$ instanton partition function from the affine–Virasoro vertex-operator-algebra identity under $\frac{\epsilon_1-\epsilon_2}{\epsilon_2}=b^2=-\frac{k+3}{k+2}$, including the required normalization of conformal blocks.
\subsection{Short English statement}
Derive the stated instanton blow-up formula from the affine--Virasoro representation identity, including matching the conformal-block normalizations.
\subsection{Sources}
\begin{itemize}
\item[\hypertarget{source:30005075:1}{S1}] ULAM.AI and the UnsolvedMath Contributors, UnsolvedMath, record 30005075 (OWR-9790367-005). Catalogue curation under CC BY 4.0; definitions and English exposition added locally.
\url{https://huggingface.co/datasets/ulamai/UnsolvedMath}.
\item[\hypertarget{source:30005075:2}{S2}] Mini-Workshop: Recent Developments in Representation Theory and Mathematical Physics (2022), Oberwolfach Reports, DOI 10.4171/owr/2022/16, pp. 817--848. Original report and mathematical context.
\url{https://doi.org/10.4171/owr/2022/16}.
\end{itemize}
\label{end:30005075}
\end{document}
